%======================================================================
\section{Track ROU --- Routing: wires, pipes, hoses (planned 27 Oct 2026)}\label{sec:rou}
%======================================================================

A new track, asked for on 27 Oct (\S\ref{sec:session-27oct}). Fusion has no routing at all; SolidWorks Routing
(electrical, piping, tubing) sets the bar: a route sub-assembly, a 3D route sketch through connection points on
the fittings, wires with a from-to list, pipes with bends, elbows and cut lengths, and the BOM by length.
FreeCAD core has nothing of it. You chose to build our own Routing workspace rather than bundle the Cables
add-on (LGPL v3, alpha, built for BIM installations on Draft's Wire and BIM's Pipe), so the routing objects share
Keystone's panels, the part-properties schema, the vault BOM and the section and interference checks. The
objects are plain \code{Part::FeaturePython} with Python proxies in \file{bruceworx\_routing.py}; the geometry
and the tables are FreeCAD-free in \file{bruceworx/routing.py} with unit tests.

\paragraph{The model.} One path, several things on it.
\begin{center}
\small
\begin{tabularx}{\linewidth}{@{}l>{\hsize=1.15\hsize}X>{\hsize=0.85\hsize}X@{}}
\toprule
Object & Properties & Shape\\
\midrule
RoutePath & points in the document frame, each optionally \emph{attached} to a vertex, hole centre, LCS or face of a
  part (it follows the part); mode Spline or Polyline-with-arcs; bend radius; minimum bend radius; read-only
  Length, Bends, BendOK & a wire: B-spline through the points, or straight runs joined by tangent arcs\\
Wire & path, gauge (AWG or mm$^2$), insulation OD, colour, signal, from/to connector and pin, slack \%, read-only
  Length and Mass & a circle swept along the path\\
Bundle & path, the wires in it, read-only OD and Length & the wires, drawn together\\
Pipe (tube) & path, OD and wall or a nominal size, material, centreline bend radius (default $2\,\text{OD}$),
  read-only CutLength and BendTable & an annulus swept along the path\\
Hose & path, ID, OD, minimum bend radius (default $6\,\text{OD}$), slack \%, read-only Length, BendOK & the same\\
\bottomrule
\end{tabularx}
\end{center}
The \emph{Routing} tab of the Design workspace has three groups: Path (Route Path, Add Point, Clip), Items
(Wire, Pipe, Hose, Bundle), Reports (From-to list, Cut list). The Route Path panel (figure~\ref{fig:routing-mock})
takes points by clicking in the 3D view, keeps them in an ordered list, shows the length and the bend check live,
and previews the wire.

\begin{figure}[htbp]
\centering
\begin{tikzpicture}[x=1mm,y=1mm,font=\sffamily\scriptsize]
  % panel
  \draw[rounded corners=2mm, fill=white, draw=black!25] (0,0) rectangle (70,78);
  \node[anchor=west, font=\sffamily\small\bfseries] at (3,73.5) {Route Path};
  \node[anchor=east, text=black!55] at (67,73.5) {Motor harness};
  \draw[black!15] (0,70) -- (70,70);
  \node[anchor=west, text=black!55] at (3,66) {Mode};
  \draw[rounded corners=0.8mm, draw=black!30] (20,63.5) rectangle (67,68.5);
  \node[anchor=west] at (22,66) {Polyline with bends}; \node[anchor=east, text=black!45] at (65,66) {$\blacktriangledown$};
  \node[anchor=west, text=black!55] at (3,59) {Bend radius};
  \draw[rounded corners=0.8mm, draw=black!30] (36,56.5) rectangle (67,61.5); \node[anchor=east] at (65,59) {12.0 mm};
  \node[anchor=west, text=black!55] at (3,52) {Minimum};
  \draw[rounded corners=0.8mm, draw=black!30] (36,49.5) rectangle (67,54.5); \node[anchor=east] at (65,52) {8.0 mm};
  \draw[black!15] (0,46) -- (70,46);
  \node[anchor=west, font=\sffamily\scriptsize\bfseries] at (3,42.5) {Points (click in the view)};
  \node[anchor=west] at (5,38) {1 \ Motor connector J1 \ \textcolor{black!50}{(face centre)}};
  \node[anchor=west] at (5,34) {2 \ Clip A \ \textcolor{black!50}{(LCS)}};
  \node[anchor=west] at (5,30) {3 \ Clip B \ \textcolor{black!50}{(LCS)}};
  \fill[bwblue!14, rounded corners=0.6mm] (3,24.2) rectangle (67,28.2);
  \node[anchor=west] at (5,26.2) {4 \ Board J3 pin 1 \ \textcolor{black!50}{(hole centre)}};
  \node[anchor=west, text=black!45] at (5,21) {$\uparrow\ \downarrow$ reorder \quad $\times$ remove};
  \draw[black!15] (0,18) -- (70,18);
  \node[anchor=west] at (3,14) {Length}; \node[anchor=east] at (67,14) {\textbf{186.4 mm}};
  \node[anchor=west] at (3,10) {Bends}; \node[anchor=east] at (67,10) {2, smallest 12.0 mm};
  \node[anchor=west] at (3,6) {Check}; \node[anchor=east, text=bwgreen] at (67,6) {\textbf{OK}};
  % the route
  \begin{scope}[xshift=80mm]
    \draw[black!20, fill=black!4] (0,10) rectangle (50,16);   % bracket
    \draw[black!20, fill=black!4] (2,16) rectangle (14,30);   % motor
    \node[text=black!50] at (8,23) {motor};
    \draw[black!20, fill=black!4] (36,16) rectangle (50,24);  % board
    \node[text=black!50] at (43,20) {board};
    % path: P1 (14,24) -> P2 (24,24) -> corner -> (24,34) -> P3 (34,34)-> corner -> (44,24)? simplified L-shapes
    \fill[bwblue] (14,24) circle (0.9); \node[anchor=south, text=bwblue] at (14,25) {1};
    \fill[bwblue] (24,24) circle (0.9); \node[anchor=south, text=bwblue] at (24,25) {2};
    \fill[bwblue] (24,36) circle (0.9); \node[anchor=south, text=bwblue] at (24,37) {3};
    \fill[bwblue] (43,36) circle (0.9); \node[anchor=south, text=bwblue] at (43,37) {4};
    \draw[bwred, very thick] (14,24) -- (20,24) arc[start angle=270, end angle=360, radius=4] -- (24,32)
         arc[start angle=180, end angle=90, radius=4] -- (43,36);
    \draw[black!40, dashed] (20,24) -- (24,24) -- (24,28);
    \node[anchor=west, text=black!55] at (25,29) {$R$};
    \draw[black!40, dashed] (24,32) -- (24,36) -- (28,36);
    \node[anchor=west, text=black!55] at (30,30) {$t=R\tan\frac{\theta}{2}$};
    \draw[black!40] (29.5,30) -- (25.5,32.5);
    \node[anchor=west, text=black!55] at (2,4) {red: the centreline, two $90^\circ$ bends of radius $R$ at points 2 and 3};
  \end{scope}
\end{tikzpicture}
\caption{The Route Path panel (left) and the path it makes (right): four picked points, straight runs joined by
tangent arcs. The corner points stay the design intent; the arcs are derived, so changing the bend radius or
moving the motor re-routes the wire.}
\label{fig:routing-mock}
\end{figure}

\paragraph{Theory: the centreline.} A route is a space curve $\mathbf r(s)$ through the picked points
$\mathbf p_0,\dots,\mathbf p_n$.
\begin{enumerate}
  \item \textbf{Spline mode} (wires, hoses): a cubic B-spline interpolating the points with centripetal
        parametrisation ($\Delta u_i=\|\mathbf p_{i+1}-\mathbf p_i\|^{1/2}$), which avoids the loops and overshoot of
        uniform parametrisation on unevenly spaced points. Its curvature
        $\kappa=\|\mathbf r'\times\mathbf r''\|/\|\mathbf r'\|^3$ is sampled along the curve; the tightest bend
        is $R_{\min}=1/\kappa_{\max}$ and the check is $R_{\min}\ge R_{\text{allowed}}$. Length is
        $L=\int\|\mathbf r'\|\,du$ by Gauss quadrature per span (OCC's \code{GCPnts\_AbscissaPoint} gives the same).
  \item \textbf{Polyline with bends} (tubes, pipes, stiff cable): straight runs $\mathbf p_{i-1}\mathbf p_i$ with a
        tangent arc of radius $R$ at each interior point. With the turn angle $\theta_i$ between consecutive runs,
        the arc starts and ends a \emph{tangent length} $t_i=R\tan(\theta_i/2)$ before and after the corner,
        and has length $R\,\theta_i$. The straight between two bends is $\ell_i-t_i-t_{i+1}$ and must be
        $\ge0$ (else the two bends overlap and the path is flagged). The developed, or cut, length is
        \[
          L_{\text{cut}}=\sum_i \ell_i-\sum_i 2t_i+\sum_i R\,\theta_i ,
        \]
        which is what a tube bender needs (a straight tube of $L_{\text{cut}}$ becomes the routed shape without
        stretch, neglecting the small bend elongation that bender tables absorb).
  \item \textbf{The bend table (LRA).} For a bender the route is described bend by bend: \emph{L}, the length fed
        from the previous bend's end tangent to this bend's start tangent (the straight above); \emph{R}, the rotation
        of the tube about its own axis between the two bend planes, the signed dihedral angle between the planes
        $(\mathbf d_{i-1},\mathbf d_i)$ and $(\mathbf d_i,\mathbf d_{i+1})$ where $\mathbf d$ are the run directions,
        $R_i=\operatorname{atan2}\big((\mathbf n_{i-1}\times\mathbf n_i)\cdot\mathbf d_i,\ \mathbf n_{i-1}\cdot\mathbf n_i\big)$
        with $\mathbf n_i=\mathbf d_i\times\mathbf d_{i+1}$ normalised; and \emph{A}, the bend angle $\theta_i$.
        A planar route has $R=0$ or $180^\circ$ throughout.
\end{enumerate}

\paragraph{Theory: wires.} The AWG gauge $n$ gives the conductor diameter
$d=0.127\cdot92^{(36-n)/39}$\,mm (so 22\,AWG is 0.644\,mm, 18\,AWG 1.024\,mm) and area $A=\pi d^2/4$; metric wire
is given by area directly (0.34, 0.5, 0.75, 1.0, 1.5, 2.5\,mm$^2$). A wire's length is the path length plus slack
(5\,\% by default, for service loops and termination), its mass $m=\rho_{\text{Cu}}\,A\,L+m'_{\text{ins}}L$
with $\rho_{\text{Cu}}=8.96$\,g/cm$^3$ and an insulation estimate from the OD. Bundles add their wires' areas
with a packing factor: $\text{OD}_{\text{bundle}}\approx1.15\sqrt{\sum \text{OD}_k^2}$ for a few wires.
Hoses are checked like spline wires against a minimum bend radius (manufacturer's figure; $6\,\text{OD}$ as the
default for a braided hydraulic hose, $3\,\text{OD}$ for silicone tubing), with slack added the same way.

\begin{epic}{ROU-01}{Route path \& bend check}{M5}{\stDone}{(new, 27 Oct)}
The RoutePath object and panel above, both modes, attachments that follow their parts, the length and the bend
check, the LRA table. \emph{Acceptance:} a path through four picked points reports its length within 0.1\,mm of
the hand calculation and flags a bend under the minimum radius; a tube bender's LRA table for a three-bend tube
matches the hand calculation.
\end{epic}

\begin{epic}{ROU-02}{Electrical wires \& harness}{M5}{\stDone}{(new, 27 Oct)}
Wires and bundles on a path with gauge, colour, signal, from-to, slack, length and mass, and the \textbf{from-to
list} (CSV: wire, signal, gauge, colour, from, to, length). \emph{Acceptance:} a three-wire motor harness lists each
wire's length and from-to in the CSV; the BOM shows it by length.
\end{epic}

\begin{epic}{ROU-03}{Pipes, tubes \& hoses}{M5}{\stDone}{(new, 27 Oct)}
Pipe/tube with OD and wall or nominal size, centreline bend radius, cut length and bend table; hose with ID/OD,
minimum bend radius, slack and the warning; the \textbf{cut list} (CSV: item, size, cut length, bend table).
\emph{Acceptance:} a three-bend tube's cut length and LRA table match the hand calculation; a hose bent under its
minimum radius is flagged.
\end{epic}

\begin{epic}{ROU-04}{Routed items in the BOM \& drawings}{M5}{\stDone}{(new, 27 Oct)}
\code{QuantityUnit} and \code{QuantityPer} in the part-properties schema (m for wires and tubes), read by the
assembly BOM and the vault BOM, which was deferred in \S\ref{sec:pdm06-plan} ``until the first harness needs it''.
Later: a flattened harness drawing (nail board) and a tube drawing with the bend table (\eref{DRW-02}).
\emph{Acceptance:} the gearbox-plus-harness BOM shows ``Wire 22\,AWG red, 0.42\,m'' as one line.
\end{epic}

\paragraph{Out of scope for now.} Auto-routing (shortest path avoiding parts). Connectors with pin maps,
fittings, harness flattening and the tube drawing were out of scope on 27 Oct; you asked for them on 30 Oct, and
they are planned in \S\ref{sec:rou-next-plan} as \eref{ROU-05}--\eref{ROU-07}.

\paragraph{Questions for you.}
\begin{enumerate}[label=R\arabic*.]
  \item \textbf{Gauge default.} AWG (your motors and boards) or mm$^2$ (metric suppliers) as the default in the
        Wire panel?\\
        \emph{Recommendation:} AWG, with mm$^2$ one click away and both printed in the from-to list.
  \item \textbf{Bend radius defaults.} Tube $2\times$OD (a typical mandrel bender), hose $6\times$OD, wire
        $4\times$OD. Keep these, or add a per-vault table (materials settings, \S\ref{sec:pdm11-plan})?\\
        \emph{Recommendation:} keep the defaults now; a per-material minimum bend radius joins the materials table
        when \eref{PDM-11} gets stock forms.
  \item \textbf{Clips.} A Clip is just an LCS on a part today. Should it be its own object with an inner diameter
        that checks the bundle fits through it?\\
        \emph{Recommendation:} yes, in \eref{ROU-04}'s second half; the LCS is enough to route with.
  \item \textbf{Wires in the BOM.} One line per wire, or one line per gauge and colour with the total length
        (what a purchasing BOM wants)?\\
        \emph{Recommendation:} per wire in the indented BOM, per gauge and colour in the flattened one.
  \item \textbf{Where routes live.} Inside the assembly file (SolidWorks makes a route sub-assembly), or as its own
        vault item linked to the assembly (a harness part number, class \emph{harness} from \S\ref{sec:pdm08-plan})?\\
        \emph{Recommendation:} its own item with a harness P/N, since it is bought and built separately.
\end{enumerate}

\InputIfFileExists{chapters/73-routing-built}{}{}
\InputIfFileExists{chapters/73b-routing-built-30oct}{}{}
%======================================================================
\subsection{BW-ROU next: connectors, fittings, drawings, the harness as a vault item (plan, 30 Oct 2026)}\label{sec:rou-next-plan}
%======================================================================

On 30 Oct you asked for all four follow-ups to the routing track (\S\ref{sec:session-30oct}). Two of them, the
Assembly BOM \emph{Length} column and the configurable bend rules, need no design and are being built now
(\S\ref{sec:rou-built-30oct}). The other three change the data model, so they come here first with their
questions. They build in this order, because each one needs the one before it:

\begin{center}
\small
\begin{tabular}{@{}lll@{}}
\toprule
Step & What & Needs\\
\midrule
1 & Connectors (pins) and fittings (ports): \eref{ROU-05} & ROU-01..04\\
2 & The harness as a vault item, the purchasing BOM by length: \eref{ROU-06} & step 1, PDM-04 vault\\
3 & Drawings: flattened harness, tube with bend table: \eref{ROU-07} & steps 1--2, DRW-01\\
\bottomrule
\end{tabular}
\end{center}

%----------------------------------------------------------------------
\subsubsection*{Step 1: connectors and fittings}

\paragraph{The idea.} A connector or fitting is an ordinary part (a vault item, often COTS) with
\textbf{named attachment frames} in it: \code{Pin1}, \code{Pin2}, \dots{} on a connector and \code{Port1},
\dots{} on a fitting. Each frame is an LCS whose origin is where the conductor or tube enters and whose $+Z$
points out along the entry direction. The frames carry a few properties:
\begin{center}
\small
\begin{tabularx}{\linewidth}{@{}lX@{}}
\toprule
Pin & name (``1'', ``A3'', ``VBAT''), accepted gauge range (e.g.\ 28--22\,AWG), terminal P/N (crimp contact),
      seal P/N, default signal\\
Port & port type (push-to-connect, barb, compression, NPT/BSPP thread), tube or hose size it takes, minimum
      straight length before the first bend\\
\bottomrule
\end{tabularx}
\end{center}
Routing then gets three behaviours: (1) a route path's first or last point picked on a pin or port
\emph{attaches} to that frame (ROU-01 already attaches to LCS origins); (2) the path \emph{leaves along the
frame's $+Z$} for the minimum straight length, so a wire does not bend at the crimp and a tube enters a fitting
straight; (3) a wire on that path fills its \code{From}/\code{FromPin} (\code{To}/\code{ToPin}) from the
connector's label and the pin's name, and the from-to list stops being typed by hand. Mating pairs (a plug and
its receptacle) record each other's P/N, so the from-to list can show a connection through a mated pair.

\paragraph{Where the geometry comes from.} Most robotics and drone connectors (JST-XH/PH/GH/SH, Molex
Micro-Fit, XT30/XT60, Dupont 2.54\,mm headers) already have STEP models in KiCad's 3D library. The KiCad
bridge (\eref{ECD-05}) already brings KiCad parts into the Standard Libraries vault, and the footprint's pads
give the pin positions and names. A connector made this way is a KiCad-library item plus the pin frames
generated from its pads (question C1). Fittings (push-to-connect, barbs, NPT elbows) have no such library.
They are made as parametric parts from a size table, as the tube sizes are in ROU-03.

\paragraph{Terminals, seals and sleeving.} A crimped wire end needs a terminal (and in sealed connectors a
seal). When a wire attaches to a pin with a terminal P/N, the terminal appears in the BOM as a counted line
under the harness (question C3). Heat-shrink and braided sleeving are cut-to-length items on a bundle and go
to the BOM by length, like wire.

%----------------------------------------------------------------------
\subsubsection*{Step 2: the harness as a vault item}

\paragraph{The problem.} You chose in R5 (\S\ref{sec:rou}) that a harness is its own vault item with a harness
P/N, because it is bought or built as one thing, with one P/N and one drawing. But it is
routed \emph{in context}: its paths attach to the motor, the board and the frame in the parent assembly. In
FreeCAD the cleanest way to keep both is the SolidWorks pattern of a \emph{virtual} component that is saved out
later:
\begin{enumerate}[nosep]
  \item a \textbf{Harness} container (an \code{App::Part} with a \code{HarnessType} marker) holds the paths,
        wires, bundles, sleeving and the connector instances, inside the assembly file, while you route;
  \item \textbf{Make vault item} saves it out as its own file under a new P/N (document type \emph{Cable
        assembly} / \emph{Tube assembly}). The parent links to it, and the path points keep their attachments
        as cross-document links (\code{App::PropertyXLink}, which FreeCAD resolves when both files are
        open, as in-context editing \eref{ASM-07} does);
  \item opened on its own, the harness keeps the last-solved point positions, so its drawing and BOM do not
        need the parent.
\end{enumerate}

\paragraph{The purchasing BOM.} Wire, tube, hose and sleeving become \textbf{stock items}: one vault item per
kind, e.g.\ ``Wire, 22\,AWG, red, PVC hook-up'', with \code{QuantityUnit}\,=\,m (question V2). Every routed
item references its stock item. The harness's own BOM then lists its connectors and terminals by count and
its stock items by length. The flattened vault BOM \emph{sums} the lengths per stock item across all harnesses
and adds a scrap allowance, the rounding to a purchase unit (a 30\,m spool) being optional:
\[
  L_{\text{buy}}=\Bigl\lceil \frac{(1+s)\sum_k n_k L_k}{L_{\text{unit}}}\Bigr\rceil L_{\text{unit}},
\]
with $n_k$ the number of harnesses of kind $k$ in the product and $L_k$ the length of that stock item in one
of them.

%----------------------------------------------------------------------
\subsubsection*{Step 3: drawings}

\paragraph{Flattened harness drawing.} The harness is a tree: connectors are leaves, breakouts are inner
nodes and bundle segments are edges with lengths. The drawing lays that tree out flat in the sheet's plane
(figure~\ref{fig:rou-flat}). It keeps the topology and labels every segment with its true length, but it is
\emph{not to scale}, as most harness drawings are (question D1). The layout is a straightforward tree drawing:
the longest path (trunk) runs horizontally, branches leave at $\pm45^\circ$ alternating sides, and a
force-directed pass separates labels that overlap. The sheet carries a connector table (connector, P/N,
mating P/N, a pin-out table per connector), the from-to wire list, the stock-item list with lengths, and
notes (lengths $\pm$ tolerance, labels, sleeving).

\begin{figure}[htbp]
\centering
\begin{tikzpicture}[x=1mm,y=1mm,font=\sffamily\scriptsize]
  \draw[line width=1.6pt, black!70] (0,20) -- (40,20) -- (85,20);
  \draw[line width=1.2pt, black!70] (40,20) -- (55,35);
  \draw[line width=1.2pt, black!70] (40,20) -- (55,5);
  \foreach \p/\n in {{(0,20)}/J1, {(85,20)}/J2, {(55,35)}/M1, {(55,5)}/S1}{
    \fill[bwblue] \p ++(-2.5,-2.5) rectangle ++(5,5);
    \node[above=2.2mm] at \p {\n};}
  \fill[black] (40,20) circle (0.9);
  \node[below=1.2mm] at (20,20) {180 $\pm$5};
  \node[below=1.2mm] at (63,20) {220 $\pm$5};
  \node[anchor=east] at (47,30) {90};
  \node[anchor=east] at (47,10) {120};
  \node[anchor=west, text=black!55] at (90,30) {J1 board, J2 battery};
  \node[anchor=west, text=black!55] at (90,25) {M1 motor, S1 sensor};
  \node[anchor=west, text=black!55] at (90,18) {segment lengths true, drawing not to scale};
  \node[anchor=west, text=black!55] at (90,13) {breakout ($\bullet$) dimensioned from J1};
\end{tikzpicture}
\caption{A flattened harness: the routed 3D tree laid out flat, trunk horizontal, branches at $45^\circ$,
every segment labelled with its routed length. Connector, pin-out and wire tables sit beside it.}
\label{fig:rou-flat}
\end{figure}

\paragraph{Tube drawing with bend table.} A tube bender works from the \textbf{YBC} (also called LRA) table:
feed length $Y$, rotation $B$ and bend angle $C$ per bend. ROU-03 already computes it. With the intersection
points $\mathbf P_0\dots\mathbf P_{n+1}$ of the polyline, the unit legs
$\mathbf d_i=(\mathbf P_{i+1}-\mathbf P_i)/\lVert\cdot\rVert$, the bend radius $R$ and $t_i=R\tan(C_i/2)$:
\[
  C_i=\arccos(\mathbf d_{i-1}\cdot\mathbf d_i),\qquad
  Y_i=\lVert\mathbf P_{i}-\mathbf P_{i-1}\rVert-t_{i-1}-t_i,\qquad
  B_i=\angle\bigl(\mathbf d_{i-2}\times\mathbf d_{i-1},\ \mathbf d_{i-1}\times\mathbf d_i\bigr),
\]
where $B_i$ is the signed dihedral angle between successive bend planes about $\mathbf d_{i-1}$ (with $t_0=0$
at the start). The developed (cut) length is $\sum Y_i+\sum R\,C_i$ plus the last straight. The drawing has an
isometric and a top view of the tube, the YBC table, the XYZ intersection-point table (what CNC benders also
accept), and the tube size, material, CLR and cut length in the title block notes. Springback and elongation
compensation depend on the bender and the material, so they are left out for now (question D3).

%----------------------------------------------------------------------
\subsubsection*{New epics}
\begin{itemize}[nosep]
  \item \eref{ROU-05} Connectors and fittings: pin and port frames, attach and leave straight, auto from-to,
        terminals in the BOM, starter library.
  \item \eref{ROU-06} Harness as a vault item: harness container, Make vault item, stock items, purchasing BOM
        by length with scrap and rounding.
  \item \eref{ROU-07} Routing drawings: flattened harness drawing with tables, tube drawing with YBC and XYZ
        tables.
\end{itemize}

\paragraph{Questions for you.}
\begin{enumerate}[label=C\arabic*.]
  \item \textbf{Connector geometry.} KiCad 3D models with pins generated from the footprint pads (through
        the ECD-05 bridge), parametric outlines we draw ourselves, or manufacturer STEP files you download?
        \emph{Recommendation:} KiCad models and pads; parametric only where KiCad has none.
  \item \textbf{Starter library.} Which families first? Proposed: JST-XH, JST-PH, JST-GH, JST-SH, Dupont
        2.54\,mm, Molex Micro-Fit 3.0, XT30, XT60, XT90, MR30, 3.5\,mm bullet (motor leads), Deutsch DT.
  \item \textbf{Terminals in the BOM} automatically from the pin's terminal P/N (one per attached wire end)?
        \emph{Recommendation:} yes.
\end{enumerate}
\begin{enumerate}[label=F\arabic*.]
  \item \textbf{Fitting families first.} Push-to-connect (4/6/8\,mm pneumatic), hose barbs, NPT/BSPP threaded
        elbows and tees, compression? \emph{Recommendation:} push-to-connect, barbs and NPT/BSPP; compression
        later.
  \item \textbf{Straight length at a pin or port.} A minimum straight length before the first bend, set on the
        pin or port (defaults: wire $2\times$OD past the terminal, tube $2\times$OD or the bender's clamp
        length), configurable like the bend rules? \emph{Recommendation:} yes.
\end{enumerate}
\begin{enumerate}[label=V\arabic*.]
  \item \textbf{Harness model.} How to reach R5's own vault item: route in context as a harness container,
        then \emph{Make vault item} saves it out (above), or each harness is its own file from the start and is
        edited in context?
        \emph{Recommendation:} the container first, saved out later.
  \item \textbf{Stock items.} Wire, tube, hose and sleeving as vault items with their own P/N per gauge,
        colour and insulation (the purchasing BOM sums by P/N), or grouped by properties without a P/N?
        \emph{Recommendation:} vault items with P/N, document type \emph{Stock material}, shared with
        \eref{PDM-11}'s stock forms.
  \item \textbf{Scrap and rounding.} A scrap allowance (default 10\,\%) and optional rounding to the purchase
        unit (spool length) per stock item? \emph{Recommendation:} yes, both per stock item, scrap default
        10\,\%.
\end{enumerate}
\begin{enumerate}[label=D\arabic*.]
  \item \textbf{Harness drawing scale.} Not to scale with every segment dimensioned (the usual practice), or
        a 1:1 formboard print (tiled PDF) as well? \emph{Recommendation:} not to scale now; the 1:1 formboard
        is a later option.
  \item \textbf{Harness drawing content.} Connector table with pin-outs, from-to list, stock-item lengths,
        labels and notes as above? Anything your builds need beyond that (wire labels/markers, test table)?
  \item \textbf{Tube drawing.} YBC and XYZ tables; springback compensation later?
        \emph{Recommendation:} yes, springback later.
\end{enumerate}

\InputIfFileExists{chapters/75-rou05-built}{}{}
\InputIfFileExists{chapters/75b-connlib-built}{}{}
%======================================================================
\subsection{BW-ROU-08 plan: flexible wires that keep their length (30 Oct 2026)}\label{sec:rou-flex-plan}
%======================================================================

\paragraph{The request.} You asked, through the session building the battery (\code{STD-02-01001-001}: two
12\,AWG leads from the pack to a T plug, four 22\,AWG balance leads to an XHP-4), for wires that behave like
real wires. Today a wire is a spline through its path points. Dragging the plug re-routes it only on release,
and the spline stretches (81\,mm became 132\,mm in the test). Your answers, already given:
\begin{enumerate}[nosep]
  \item the wire \textbf{keeps its length} (its cut length) and never bends tighter than its minimum radius;
        the ends stay clamped to the pin frames and leave them straight, as now;
  \item it moves \textbf{live during an Assembly drag}, solved again at each drag step from the last shape;
  \item pulled taut, \textbf{the drag stops} at the wire's length, like the joint-limit drag stop of
        \S\ref{sec:joints} (no warn and continue);
  \item \textbf{gravity sag} is in;
  \item \textbf{collisions with parts} are in, against simplified shapes for speed;
  \item this is the \textbf{default for every wire}, not a mode.
\end{enumerate}
The method borrows the \emph{ideas} of Blender's hair and cloth tools, not their code (Blender is GPL): position-based
dynamics for the rope and discrete elastic rods for the stiffness. Everything is numpy; nothing new is bundled.

%----------------------------------------------------------------------
\subsubsection*{Theory}

\paragraph{The rope.} A wire of cut length $L$ is $N$ points $\mathbf x_0\dots\mathbf x_{N-1}$ with a rest
segment length $\ell_0=L/(N-1)$. The first and last $k$ points sit on the leads: they are fixed on the pin
frames' $+Z$ axes for the straight lengths of ROU-05, so the wire leaves each terminal straight. The rest move.
The shape is the \emph{resting} configuration under three kinds of constraint and gravity. It is not a dynamic
simulation: we want where the wire settles, not how it swings.

\paragraph{XPBD (Macklin, M\"uller, Chentanez 2016).} Extended position-based dynamics moves points directly to
satisfy constraints $C_j(\mathbf x)=0$ (or $\ge0$), each with a compliance $\alpha_j$ (the inverse of a
stiffness). One step of length $\Delta t$:
\[
  \tilde{\mathbf x}_i=\mathbf x_i+\Delta t\,\mathbf v_i+\Delta t^2\,\mathbf g,\qquad
  \Delta\lambda_j=\frac{-C_j-\tilde\alpha_j\lambda_j}{\sum_i w_i\lVert\nabla_iC_j\rVert^2+\tilde\alpha_j},\qquad
  \Delta\mathbf x_i=w_i\,\nabla_iC_j\,\Delta\lambda_j,
\]
with $\tilde\alpha_j=\alpha_j/\Delta t^2$, inverse masses $w_i$ (zero for the clamped points), and the
Lagrange multipliers $\lambda_j$ accumulated over the iterations. With heavy velocity damping and the previous
shape as the start, a few steps reach rest. That is what lets a drag step cost milliseconds, because each step
starts from where the wire already is.

\paragraph{Constraints.}
\begin{itemize}
  \item \textbf{Length} (inextensible, $\alpha=0$): $C=\lVert\mathbf x_{i+1}-\mathbf x_i\rVert-\ell_0$ per segment.
        The wire neither stretches nor shrinks; its length is $L$ by construction, not by checking.
  \item \textbf{Bending} (elastic): the turning angle $\theta_i$ at point $i$ between consecutive segments.
        The discrete elastic rod (Bergou et al.\ 2008) gives the bending energy
        $E_b=\sum_i \frac{EI}{\ell_0}\bigl(2\tan\frac{\theta_i}{2}\bigr)^2$, so the compliance is
        $\alpha_b=\ell_0/EI$. Twist is left out: a round wire's twist barely changes its centreline.
  \item \textbf{Minimum bend radius} (a limit, not a spring): the local radius $R_i=\ell_0/(2\sin\frac{\theta_i}{2})$
        must stay at or above $R_{\min}$ from the wire's bend rule, i.e.\ $\theta_i\le\theta_{\max}=2\arcsin
        \frac{\ell_0}{2R_{\min}}$. It is active only when violated. For 12\,AWG at $4\times$OD ($R_{\min}=14.4$\,mm)
        and $\ell_0=6$\,mm, $\theta_{\max}=24^\circ$.
  \item \textbf{Collision} (inequality): for each point and each nearby part proxy with signed distance $d(\mathbf x)$,
        $C=d(\mathbf x_i)-r\ge0$ with $r$ the wire radius plus a clearance. A violated point is pushed out along
        $\nabla d$.
\end{itemize}

\paragraph{Stiffness from the gauge.} A stranded conductor bends with its strands sliding, so its stiffness is
close to the sum of the strands', $EI_c\approx n_s E_{\text{Cu}}\pi d_s^4/64$, a lower bound (friction adds
some). The insulation is a tube, $EI_i=E_i\pi(D^4-d^4)/64$. For 12\,AWG (65 strands of 0.254\,mm, PVC OD
3.6\,mm): $EI_c\approx1554$, $EI_i\approx137$, $EI\approx1.7\times10^3$\,N\,mm$^2$. A table per insulation type
(PVC, silicone, PTFE, XLPE, matching the ROU bend-rule presets) holds $E_i$ and density, and a per-wire factor
calibrates it against a hang test (question W9).

\paragraph{Gravity.} The weight per length $q=(\rho_{\text{Cu}}A_c+\rho_iA_i)g$ is
$3.7\times10^{-4}$\,N/mm for 12\,AWG. Gravity points along $-Z$ of the top assembly unless the assembly says
otherwise (question W7).

\paragraph{Taut and the drag stop.} An inextensible rope whose end points move apart beyond its length has no
solution. The solver detects this as a residual: the length constraints cannot all be met, and the largest
segment strain stays above $10^{-3}$ after the iterations. The quick test first is the lower bound
$\lVert\mathbf a-\mathbf b\rVert+\ell_a+\ell_b\le L$ between the lead tips $\mathbf a,\mathbf b$ (with obstacles
the shortest path is longer, so the residual decides). The drag step is then cut back by bisection, as the joint
limits already are (\S\ref{sec:joints}): the part stops where the wire just reaches its length.

\begin{figure}[htbp]
\centering
\begin{tikzpicture}[x=1mm,y=1mm,font=\sffamily\scriptsize]
  \fill[black!8] (0,10) rectangle (12,34); \node at (6,36) {pack};
  \fill[black!8] (108,16) rectangle (120,40); \node at (114,42) {plug};
  \fill[bworange!25] (48,0) rectangle (70,12); \node at (59,6) {part (proxy box)};
  \draw[black!40, dashed] (12,22) -- (20,22); \draw[black!40, dashed] (100,28) -- (108,28);
  \draw[bwblue, line width=1.1pt] plot[smooth, tension=0.6] coordinates {(12,22) (20,22) (32,18) (46,15) (59,14.5) (72,16) (86,21) (100,28) (108,28)};
  \foreach \p in {(12,22),(16,22),(20,22),(108,28),(104,28),(100,28)} \fill[black] \p circle (0.8);
  \foreach \p in {(26,20),(32,18),(39,16.3),(46,15),(52.5,14.6),(59,14.5),(65.5,15),(72,16),(79,18.3),(86,21),(93,24.5)} \draw[bwblue, fill=white] \p circle (0.8);
  \draw[->, black!50] (5,4) -- (5,-2) node[right] {$\mathbf g$};
  \node[anchor=west, text=black!55] at (2,46) {$\bullet$ clamped on the lead (straight length of the pin)};
  \node[anchor=west, text=black!55] at (2,50) {$\circ$ free points: length, bending, bend limit, collision, gravity};
  \node[anchor=west, text=black!55] at (60,50) {the part proxy pushes points out by $r$};
\end{tikzpicture}
\caption{A wire as a rope of points. The leads are clamped along the pin frames; the free points settle under
gravity while every segment keeps $\ell_0$, every bend stays above $R_{\min}$, and no point enters a part proxy.}
\label{fig:rou-flex}
\end{figure}

%----------------------------------------------------------------------
\subsubsection*{How it fits Keystone}

\begin{itemize}
  \item \textbf{The wire keeps a state.} A new hidden property stores the rope's points. A recompute that moves a
        pin (an edited joint, a moved part) solves from that state, so a wire never jumps to a new shape.
  \item \textbf{During a drag.} Assembly drags go through \code{preDrag}/\code{doDragStep}/\code{postDrag}
        (commit 1a15d0d1). A small \code{[BW]} hook in \code{AssemblyObject::doDragStep} calls the registered
        Python \emph{drag checks} after the solver proposes placements and before they are applied. The rope
        module is one of them: it solves the wires attached to the moved parts, returns ``ok'' or ``taut'', and
        a taut answer sends the step into the existing bisection. While dragging, a wire is drawn as a light
        tube (a Coin line set with the wire's width and colour), because sweeping an OCC solid 30 times a second is
        too slow; the real solid is built once on release.
  \item \textbf{Collision proxies.} Each part near a wire (bounding boxes within the wire's reach) gets a proxy:
        its oriented bounding box by default, or a convex hull or decimated mesh for parts where the box is too coarse
        (question W6). Signed distances to boxes are exact and cheap; to a mesh, through a small BVH in numpy.
  \item \textbf{Resolution.} About one point per 1.5\,OD, clamped to 20--200 points (question W4); 40 for the
        battery's 81\,mm lead.
  \item \textbf{Flexible sub-assemblies.} An \code{AssemblyLink} with \code{Rigid=False} copies only links and
        assembly links into the parent (\code{synchronizeComponents}). Wires and route paths stay behind in the
        sub-assembly's document and do not follow parts the parent moves. The fix: the copy step also copies
        routed items, re-pointing their attachments to the copied links (C++, \code{[BW]}).
\end{itemize}

\paragraph{Verification.}
\begin{center}
\small
\begin{tabularx}{\linewidth}{@{}lX@{}}
\toprule
Catenary & $EI\to0$, ends 200\,mm apart at the same height, $L=250$\,mm: the shape is a catenary with
  $a=84.55$\,mm and sag $a(\cosh\frac{d}{2a}-1)=66.36$\,mm; the test accepts 1\,\%.\\
Cantilever stub & 12\,AWG, 100\,mm clamped horizontally at one end, free at the other: tip sag
  $qL^4/8EI=2.76$\,mm (small deflection); accepts 5\,\%.\\
Length & any drag sequence: $\sum\lVert\mathbf x_{i+1}-\mathbf x_i\rVert=L$ within $10^{-3}$.\\
Bend limit & a pin turned back on itself: no turning angle above $\theta_{\max}$.\\
Drag stop & the T plug dragged away from the pack stops where the 12\,AWG leads are taut, within one bisection
  step; dragging back frees it.\\
Collision & a wire over a box sags onto it and no point is inside the box minus $10^{-2}$\,mm.\\
Battery & \code{STD-02-01001-001} (on a vault copy): drag P1 through an arc; the leads keep 81\,mm, update live
  ($\ge20$ frames/s), stop when taut.\\
\bottomrule
\end{tabularx}
\end{center}

\paragraph{New epic.} \eref{ROU-08} Flexible wires: rope solver as the default wire behaviour, live in drags, drag
stop when taut, gravity, collision proxies, flexible sub-assemblies carrying routed items.

\paragraph{Questions for you.}
\begin{enumerate}[label=W\arabic*.]
  \item \textbf{Hand-placed path points.} Under the rope, what is a point you picked? \emph{Recommendation:} a
        point picked on a part (a clip, an LCS) becomes a \emph{clip} the wire passes through, free to turn; a point
        in free space is only a \emph{guide} for the first shape and then forgotten.
  \item \textbf{Migrating existing wires.} Each wire's cut length becomes its current length (with its slack), so
        nothing changes shape until something moves? \emph{Recommendation:} yes, and the slack field then edits the
        cut length directly.
  \item \textbf{Bundles, hoses, pipes.} \emph{Recommendation:} hoses follow (same rope, hose stiffness); a bundle is
        one rope with the combined OD and stiffness, branches at breakouts later; pipes and tubes stay rigid (they are
        bent, not flexible).
  \item \textbf{Resolution.} One point per 1.5\,OD, 20--200 points? \emph{Recommendation:} yes.
  \item \textbf{Live display.} A light tube while dragging, the real solid on release? \emph{Recommendation:} yes.
  \item \textbf{Collision proxies.} Oriented box per part by default, convex hull or mesh chosen per part; wire to
        wire contact later? \emph{Recommendation:} yes.
  \item \textbf{Gravity.} $-Z$ of the top assembly, settable per assembly (a drone shown on its side)?
        \emph{Recommendation:} yes.
  \item \textbf{Outside drags.} A recompute that moves a pin (joint edit, typed placement) solves the rope from its
        last shape, and a taut wire there turns red with a badge rather than refusing the edit?
        \emph{Recommendation:} yes. Only drags can be stopped.
  \item \textbf{Stiffness data.} Table per insulation type plus a per-wire calibration factor from a hang test (clamp
        100\,mm, measure the tip sag)? \emph{Recommendation:} yes.
  \item \textbf{Bend limit while dragging.} If a drag would force a bend tighter than the minimum (the pin turned hard),
        stop the drag there too, or let the wire turn red? \emph{Recommendation:} stop the drag, like taut.
  \item \textbf{Flexible sub-assemblies.} Copy routed items into the parent with the links (above)?
        \emph{Recommendation:} yes.
\end{enumerate}

\paragraph{Your answers (30 Oct).} All recommendations W1--W11 accepted, with three additions:
\begin{itemize}[nosep]
  \item \textbf{Any length.} A wire or hose can be made at any cut length you type, longer than its route (it
        sags, loops or lies on the parts below) or as short as its leads allow (then it is taut, and red outside a
        drag). The cut length is the wire's own property, not derived from the path; the path's points only give
        the first layout (W1).
  \item \textbf{Existing wires do not extend}: they keep their present length as their cut length (W2), and no
        wire ever stretches.
  \item \textbf{Gravity is always on}, and its direction is chosen per assembly along an axis ($\pm X$, $\pm Y$,
        $\pm Z$ of the assembly, default $-Z$), W7.
\end{itemize}
Also accepted: Insert Fastener files new bolts under a \emph{Hardware} category (\S\ref{sec:categories-built}).

\InputIfFileExists{chapters/76b-rou-flex-built}{}{}
\InputIfFileExists{chapters/76c-flexsub-built}{}{}
